% =====================================================================
\section{Conclusion}
% =====================================================================
Open, locally-servable LLMs are now strong enough to serve as annotators on objective and verifiable tasks, where
they meet or exceed established proprietary-judge correlations. But across \Nstudies{} reviewed studies and a
\STHRn{}-judgment meta-evaluation on six datasets with real human labels, reliability is bounded by the same thing
for every model: the \emph{human-human agreement ceiling}, and it degrades on exactly the items where humans
themselves disagree. Scaling the judge helps only modestly, and contamination does not explain the pattern. The
practical consequence is that an LLM-as-a-judge should be reported the way we report any measurement instrument,
against a human reference, with chance-corrected and disagreement-aware metrics, with uncertainty, and with a stated
operating range, rather than trusted as an oracle. The right question is never ``can a machine judge?'' but ``can
this open model agree with humans, on this construct, as well as humans agree with each other?''
